%! Absolute Value Functions \& Transformations — Unit 2, Lesson 2.3 — Warm-Up (Answer Key)
\documentclass[10pt]{article}
\usepackage{algebra2-article}
\usepackage{algebra2-key}   % pulls in -boxes; do NOT also load -boxes
\newcommand{\TallMath}[1]{$\displaystyle #1\rule[-1.4em]{0pt}{3.2em}$}
% Coordinate grid: {xmin}{xmax}{ymin}{ymax}
\newcommand{\lgrid}[4]{%
  \draw[step=1, linegray!40, very thin] (#1,#3) grid (#2,#4);
  \draw[->, semithick, charcoal] (#1,0)--(#2,0) node[below right, font=\scriptsize]{$x$};
  \draw[->, semithick, charcoal] (0,#3)--(0,#4) node[left, font=\scriptsize]{$y$};}

\begin{document}

\pageheader{Unit 2, Lesson 2.3}{Warm-Up --- Answer Key}

\vspace{0.10in}

\begin{notesbox}{Getting Ready: Distance, Transformations, and the V}
\small Today an absolute value becomes a \emph{graph}. These three quick reviews are the tools you'll
need to read it.

\vspace{4pt}
\textbf{1. Build the table for $y=|x|$.} Evaluate each absolute value.
\begin{center}
\renewcommand{\arraystretch}{1.25}
\begin{tabular}{|c|c|c|c|c|c|}
\hline
$x$ & $-2$ & $-1$ & $0$ & $1$ & $2$ \\ \hline
$y=|x|$ & ${\color{keyred}\mathbf{2}}$ & ${\color{keyred}\mathbf{1}}$ & ${\color{keyred}\mathbf{0}}$ & ${\color{keyred}\mathbf{1}}$ & ${\color{keyred}\mathbf{2}}$ \\ \hline
\end{tabular}
\end{center}
The two sides of the table are \ans{the same} (the same / different), and the smallest
$y$-value is \ans{$0$}, at $x=\ans{$0$}$.

\vspace{6pt}
\textbf{2. Transformations of the parent (from Lesson 2.1).} Fill each blank with a direction or word.
\begin{itemize}[leftmargin=1.5em, itemsep=2pt]
  \item Compared to $y=x$, the graph of $y=x+3$ shifts \ans{up} by $3$ (up / down).
  \item Multiplying a parent function by a \emph{negative} number \ans{reflects} it over the $x$-axis.
\end{itemize}

\vspace{6pt}
\textbf{3. Read the V.} The graph of the parent $y=|x|$ is shown.
\begin{center}
\begin{tikzpicture}[scale=0.5]
  \lgrid{-4}{4}{-1}{5}
  \draw[very thick, forest] (-4,4)--(0,0)--(4,4);
  \fill[charcoal] (0,0) circle (3.5pt);
  \draw[navy, dashed, semithick] (0,-1)--(0,5);
\end{tikzpicture}
\end{center}
The lowest point (the \emph{vertex}) is at $({\color{keyred}\mathbf{0}},\,{\color{keyred}\mathbf{0}})$,
and the graph is symmetric about the vertical line $x=\ans{$0$}$.
\end{notesbox}

\end{document}
